% IEEE-formatted preliminary article (conference)
\documentclass[conference]{IEEEtran}
\usepackage[utf8]{inputenc}
\usepackage[T1]{fontenc}
\usepackage[english]{babel}
\usepackage{csquotes}
\usepackage{url}
\usepackage{graphicx}
\usepackage{microtype}
\usepackage{adjustbox}
\usepackage{rotating}
\usepackage{booktabs}
\usepackage{multirow}
\usepackage{hyperref}
\usepackage{amsmath}
\usepackage{amssymb}
\usepackage{tikz}
\usepackage{xcolor}
\usepackage{listings}
% babel (english) redefine estos nombres al iniciar el documento
\addto\captionsenglish{%
  \renewcommand{\figurename}{Figura}%
  \renewcommand{\tablename}{Tabla}%
}
\renewcommand{\lstlistingname}{Código}

% Estilo para bloques de código Python
\lstset{
  language=Python,
  basicstyle=\ttfamily\footnotesize,
  keywordstyle=\color{blue!70!black}\bfseries,
  commentstyle=\color{green!45!black}\itshape,
  stringstyle=\color{red!65!black},
  numbers=left,
  numberstyle=\tiny\color{gray},
  numbersep=6pt,
  frame=single,
  rulecolor=\color{gray!60},
  backgroundcolor=\color{gray!8},
  breaklines=true,
  showstringspaces=false,
  upquote=true,
  tabsize=4,
  captionpos=b,
  % Permite tildes y ñ dentro del código (p. ej. en comentarios)
  literate={á}{{\'a}}1 {é}{{\'e}}1 {í}{{\'i}}1 {ó}{{\'o}}1 {ú}{{\'u}}1
           {Á}{{\'A}}1 {É}{{\'E}}1 {Í}{{\'I}}1 {Ó}{{\'O}}1 {Ú}{{\'U}}1
           {ñ}{{\~n}}1 {Ñ}{{\~N}}1
}
\usetikzlibrary{positioning, fit, arrows.meta, calc, patterns, decorations.pathreplacing}

% --- Utilidades del borrador ---
% Valor que se obtiene al ejecutar los scripts de Python
\newcommand{\pendiente}[1]{\textcolor{red}{\textbf{[#1]}}}
% Incluye la figura si ya existe; si no, deja un recuadro indicándolo
\newcommand{\figuraopc}[2]{%
  \IfFileExists{figuras/#1}%
    {\includegraphics[width=\columnwidth]{figuras/#1}}%
    {\fbox{\parbox{0.9\columnwidth}{\centering\vspace{1.5em}%
      \textcolor{red}{Figura pendiente: ejecutar \texttt{#2}}\vspace{1.5em}}}}}

\title{Tarea 4 - Señales y Sistemas en Tiempo Discreto 2}
\author{\IEEEauthorblockN{Zúñiga Ramírez Jose Eduardo}
\IEEEauthorblockA{Escuela de Ingeniería Electrónica \\
Instituto Tecnológico de Costa Rica \\
\{je.zuniga\}@estudiantec.cr}
}
\date{}

\begin{document}
\maketitle

% Loosen line breaking to avoid overfull boxes in two-column layout
\sloppy


\section{Pregunta 1}
Para calcular los polos y los ceros, se debe empezar por multiplicar $z^{k}$ en el numerador y denominador 
para eliminar el factor negativo del exponente y para poder encontrar el cero que está oculto. Haciendo la 
operación anterior se obtiene la ecuación \ref{eq:polinomio} y se iguala a cero tanto el numerador como el 
denominador para obtener los ceros como polos respectivamente de la función.

\begin{equation}
H(z) = \frac{(1-\alpha)\,z^{k}}{z^{k} - \alpha}
\label{eq:polinomio}
\end{equation}

\subsection{Ceros}
Para el cálculo de los ceros se tiene: $(1-\alpha)z^k = 0$, donde para que la expresión sea cero, $z^{k}$ 
debe ser cero, por lo que la expresión tiene un cero en $z = 0$ de orden $k=10$.

\subsection{Polos}
Para el caso de los polos se tiene $z^{k} - \alpha = 0$, despejando se tiene que  $z^{k} = \alpha$. Para 
obtener los polos se debe calcular la raiz $k$-ésima de $\alpha$. La expresión se puede reescribir con 
coordenadas polares para obtener todas las raices posibles de una manera más simple, teniendo la siguiente 
igualdad: $z = r\,e^{j\theta}$. Con lo anterior se pueden reescribir las expresiones como $z^k = \left(r e^{j\theta}\right)^k = r^k\, e^{jk\theta}$ 
y como $\alpha = \alpha\, e^{j2\pi m}$. Se procede a reescribir la igualdad de $z^{k} = \alpha$ con las 
nuevas expresiones y se obtiene el resultado de la ecuación \ref{eq:polos}.

\begin{equation}
r^k\, e^{jk\theta} = \alpha\, e^{j2\pi m}
\label{eq:polos}
\end{equation}

Se iguala cada parte por separado y se despeja el radio ($r$) y el angulo ($\theta$). Para el caso del radio 
se tiene $r^k = \alpha$ y despejando $r = \alpha^{1/k}$. Para el ángulo $k\theta = 2\pi m$ y despejando 
$\theta_m = \dfrac{2\pi m}{k}$. Con los resultados anteriores se pueden obtener los números en coordenadas 
polares (\ref{eq:resultado}) que corresponden a las distintas raíces de la expresión y por ende, los polos. 

\begin{equation}
polos = \alpha^{1/k}\, e^{\,j\frac{2\pi m}{k}}, \qquad m = 0, 1, \dots, k-1
\label{eq:resultado}
\end{equation}



\section{Pregunta 3}
asd




\section{Pregunta 4}
asd




\end{document}
